\begin{abstract}
Network intrusion detectors are evaluated on random splits and reach near-perfect scores. In deployment, they
meet drifting traffic and attack classes that did not exist at training time, while analysts can label only a
small fraction of the flows, and only after a delay. Chronological replays of
CIC-IDS2017, CSE-CIC-IDS2018 and UNSW-NB15 give every method the same label budget and delay, and we score them
with metrics that stay comparable when methods learn different classes. We propose \nova{}, an adaptation loop
that rejects unfamiliar flows with a class-conditional $k$-nearest-neighbour detector and recalibrates it for
newly learned classes. It spends the label budget on novel, diverse traffic, propagates analyst labels to
agreeing neighbours, keeps rare attacks in a class-balanced memory, and deploys an update or a rollback only
after a hold-out and a forgetting check. We compare \nova{} with eight methods, including re-implementations of
INSOMNIA, NetGuard and class-balanced replay. At a 10\% budget it matches the strongest
baseline on open-world macro-F1 on the CIC datasets (0.752 against 0.755 and 0.738 against 0.730) and flags significantly fewer benign flows
on all three datasets. With 1\% of CIC-IDS2017 labelled, it leads by 8.5 points. On UNSW-NB15, stale
rejection thresholds push the replay-based methods far below the never-updated model; an ablation shows that
\nova{}'s per-class recalibration avoids this. We also report where it fails: short-lived attacks that resemble benign
traffic, a neural classifier, and one dataset at the smallest budget. Code,
configurations and per-flow predictions are released.
\end{abstract}

\begin{keyword}
network intrusion detection \sep concept drift \sep open-set recognition \sep active learning \sep
continual learning \sep evaluation methodology
\end{keyword}
